\documentclass[11pt,a4paper]{article}
\usepackage[margin=2.4cm]{geometry}
\usepackage{graphicx,booktabs,amsmath,amssymb,natbib,hyperref,xcolor}
\usepackage[T1]{fontenc}
\hypersetup{colorlinks=true,linkcolor=blue!55!black,citecolor=blue!55!black,urlcolor=blue!55!black}
\graphicspath{{../plots/}}
\newcommand{\masyr}{mas\,yr$^{-1}$}
\newcommand{\kms}{km\,s$^{-1}$}
\newcommand{\ocen}{$\omega$\,Cen}
\newcommand{\Msun}{M_\odot}
\title{\ocen: the mass-modelling procedure}
\author{oCen\_dm project notes}
\date{2026 September 21}

\usepackage{microtype}
\usepackage{adjustbox}
\setlength{\emergencystretch}{2em}
\begin{document}
\maketitle

\begin{abstract}
The companion note describes how the kinematic data were prepared. This one describes what is
done with them: the mass model, the dynamical solver, the likelihood, the priors, and the
inference. The question is whether \ocen\ requires mass beyond its stars, remnants and any
central black hole, and it is posed as a comparison of Bayesian evidence between two model
families that differ by one component. Everything here is deliberately low-dimensional: the
data are 89 binned points over roughly 0.1--55\,pc. Three isotropic fits without instrument
scales are complete. They have no persuasive halo preference and all leave structured
residuals, so the posteriors are conditional baselines. This note describes the implemented
machinery and distinguishes it from the proposed anisotropy extension.
\end{abstract}

\section{The question, and how it is posed}

Two families are fitted to the same observations and compared by evidence.

\begin{description}
\item[K1, no dark matter] stars, a compact remnant population, and a central point mass.
\item[K2, dark matter] the same three components plus a truncated generalized NFW halo.
\end{description}

K1 is the zero-halo limit of the K2 equations. Zero is outside the adopted log-uniform
K2 mass prior, so it is not a sampled nested point. The Bayes factor
$\ln Z_{K2}-\ln Z_{K1}$ compares the separately specified families and their priors. Two K2 variants are run separately,
a cored halo ($\gamma=0$) and a cuspy one ($\gamma=1$), rather than fitting $\gamma$: making
the inner slope free gives the halo shape freedom the data do not support, and the two
limiting cases bracket it.

\section{The mass model}

The potential is spherical and the components add:
\begin{equation}
M(<r) = M_\star(<r) + M_{\rm rem}(<r) + M_\bullet + M_{\rm DM}(<r).
\end{equation}

\paragraph{Stars.} A multi-Gaussian expansion, deprojected analytically under spherical
symmetry, scaled by a single free total mass $M_\star$. Using an MGE rather than a Plummer
sphere matters because the tracer density appears in the Jeans equation directly, so its
shape is not a free choice.

The default profile is not the surface brightness alone. It is a \emph{composite}: HST star
counts inside $25''$ spliced onto the Trager surface-brightness profile outside, because the
light profile is unreliable in the crowded core where a handful of bright giants dominate it.
The Trager-only alternative is selectable. The same shape supplies both the tracer density
$\nu(r)$ and the stellar mass density, which assumes a radius-independent mass-to-light
ratio; mass segregation makes that an approximation.

\paragraph{Remnants.} A Plummer sphere with free mass $M_{\rm rem}$ and scale $a_{\rm rem}$.
It is a separate component rather than an inflation of the stellar mass-to-light ratio so that
its mass appears under its own name in every posterior and can be confronted with
multi-mass predictions later. The parametrisation is phenomenological: the point is to let a
centrally concentrated dark component of \emph{known} astrophysical origin compete with the
halo, so that a halo detection cannot be manufactured by forbidding one.

\paragraph{Central point mass.} $M_\bullet$, an intermediate-mass black hole or any
unresolved central mass.

\paragraph{Dark halo (K2 only).} A truncated generalized NFW profile,
\begin{equation}
\rho(r) = \rho_s\left(\frac{r}{r_s}\right)^{-\gamma}
          \left(1+\frac{r}{r_s}\right)^{\gamma-3} T(r), \qquad
T(r) = \left[1+\left(\frac{r}{r_t}\right)^{2}\right]^{-2}.
\end{equation}
The taper falls as $r^{-4}$ outside $r_t$, giving an outer density proportional to
$r^{-7}$ and finite total mass. The fitted family fixes $r_t=1000$\,pc, outside the
measured radii. This is a phenomenological taper, not a tidal boundary derived from the
cluster's orbit; it should not be confused with the much smaller radii in the survival
experiments.

The halo is parametrised by $M_{\rm DM}(<100\,\mathrm{pc})$ and $r_s$ rather than by $\rho_s$
and $r_s$. The normalisation is solved from the enclosed mass at 100\,pc. Both this mass and $r_s$
have log-uniform priors. The 100-pc mass remains an extrapolation beyond the outermost
measured radius; the scale radius is poorly constrained in the baseline runs.

\section{Anisotropy}

The implemented choices are isotropic, constant $\beta$, or the monotonic form
\begin{equation}
\beta(r) = \beta_0 + (\beta_\infty-\beta_0)\,\frac{r^2}{r^2+r_\beta^2}.
\end{equation}
The last has three parameters and a closed-form Jeans integrating factor. It is the
logistic form passed to JamPy with $\alpha=2$, although the two engines still use different
numerical approximations. It can cross zero once; it cannot have an intermediate maximum.
A turnover extension is motivated by the residuals and remains unimplemented.

For a cored tracer in a central point-mass potential, the necessary condition from
\citet{anevans2006} is $\beta(0)\le-1/2$. The varying family imposes
$\beta_0\in[-1,-1/2]$. This is necessary, not sufficient, for a non-negative distribution
function. The condition also applies to constant anisotropy: the isotropic baseline and
constant-$\beta$ prior $[-1,0.5]$ deliberately allow moment models that violate it. Their
purpose is diagnostic, and their central-mass estimates must be interpreted accordingly.

\section{The dynamical solver}

Given $\nu(r)$ from the tracer MGE, $M(<r)$ from the mass model and $\beta(r)$, the spherical
Jeans equation
\begin{equation}
\frac{\mathrm{d}(\nu\sigma_r^2)}{\mathrm{d}r} + \frac{2\beta\nu\sigma_r^2}{r}
= -\frac{\nu\,G M(<r)}{r^2}
\end{equation}
integrates with the factor $g(r)=\exp\left(2\int\beta/t\,\mathrm{d}t\right)$ to
\begin{equation}
\nu\sigma_r^2(r) = \frac{1}{g(r)}\int_r^\infty g(s)\,\nu(s)\,\frac{G M(<s)}{s^2}\,\mathrm{d}s ,
\end{equation}
and the projected second moments follow \citep{binneymamon1982}:
\begin{align}
\Sigma\sigma_{\rm los}^2(R) &= 2\int_R^\infty \left[1-\beta\frac{R^2}{r^2}\right]
   \nu\sigma_r^2\,\frac{r\,\mathrm{d}r}{\sqrt{r^2-R^2}},\\
\Sigma\sigma_{\rm pmR}^2(R) &= 2\int_R^\infty \left[1-\beta+\beta\frac{R^2}{r^2}\right]
   \nu\sigma_r^2\,\frac{r\,\mathrm{d}r}{\sqrt{r^2-R^2}},\\
\Sigma\sigma_{\rm pmT}^2(R) &= 2\int_R^\infty \left[1-\beta\right]
   \nu\sigma_r^2\,\frac{r\,\mathrm{d}r}{\sqrt{r^2-R^2}} .
\end{align}
All three projections and $\Sigma(R)$ share their quadrature nodes, so one model evaluation
delivers the line-of-sight, radial and tangential predictions together.

\paragraph{Numerics.} The substitution $r=R\cosh u$ removes the integrable singularity at
$r=R$ exactly, since $r\,\mathrm{d}r/\sqrt{r^2-R^2}=r\,\mathrm{d}u$, leaving smooth integrals
that fixed Gauss--Legendre nodes handle to high accuracy. $\nu\sigma_r^2$ is tabulated once per
model on a 600-point logarithmic grid and splined in the log. A full set of projected profiles
costs a few milliseconds. The completed rung-0 runs required 3.2--14.3 million
likelihood calls and 4.0--20.6 hours; sampling efficiency is a practical limitation.

One defect worth recording: the outer integral was originally evaluated as
$\mathrm{cumulative}[-1]-\mathrm{cumulative}$, which loses all significance once the remaining
integral is small, and $\sigma_r$ collapsed to zero around three MGE scale lengths. It is now
accumulated in reverse. The runs made before that fix are quarantined.

\section{The likelihood}

Each dataset is a set of binned dispersions with asymmetric errors. The log-likelihood is a
split normal,
\begin{equation}
\ln \mathcal{L} = \sum_i -\frac{1}{2}\left(\frac{m_i-d_i}{\varepsilon_i}\right)^2
                  + \ln\frac{\sqrt{2/\pi}}{\varepsilon_i^{-}+\varepsilon_i^{+}}, \qquad
\varepsilon_i = \begin{cases}\varepsilon_i^{+} & m_i > d_i\\ \varepsilon_i^{-} & \text{else,}\end{cases}
\end{equation}
summed over datasets, which are treated as independent. The normalisation is \emph{shared}
between the two sides, which is what makes the density continuous at $m_i=d_i$ and integrate
to one; giving each side its own $-\ln\varepsilon$, as this did until 2026-09-20, produced a
jump of $\varepsilon^{+}/\varepsilon^{-}$ across every datum.

Three things happen between the model and the datum.

\paragraph{Radial averaging.} For the datasets we measured ourselves --- the four HST and
\emph{Gaia} profiles --- the model is averaged over the radii of the stars actually measured
in each bin, eight equal-count quantiles stored with the data, rather than over a complete
annulus with $\Sigma(R)R$ weighting. Where the selection's coverage changes inside a bin the
two differ: HST's $300$--$340''$ sample has a star-weighted mean radius of $313.2''$ against
the annulus midpoint of $320''$, worth up to $1.17\sigma$. The 29 \textbf{MUSE} bins are
published, carry no such information, and are still averaged over complete annuli.

\paragraph{Streaming.} The model predicts the full second moment; the data are dispersions
about a rotating mean. Each dataset that needs it carries $\langle\bar v^2\rangle$, and the
comparison is with $\sqrt{\sigma_{\rm model}^2-\langle\bar v^2\rangle}$. The uncertainty on
that term is propagated into the datum's error for the HST and \emph{Gaia} profiles.
\textbf{It is not for MUSE}, whose rotation and dispersion were fitted jointly by the survey
and would need their covariance rather than an independent error: at the innermost bin the
first-order rotation contribution is about $1.22$\,\kms\ against quoted errors of
$1.90/2.52$\,\kms, so this is not negligible and is an outstanding item.

\paragraph{Instrument scales.} Optional multiplicative nuisances $s_{\rm inst}$ on the model
prediction, one per instrument, with uniform priors. They absorb a constant calibration
offset. They are also strongly degenerate with the halo mass, which is why the adopted family comparisons use \texttt{--no-scales}. A separate scale-freeing diagnostic
may be run after the sequence, but its evidence is not compared. The bare CLI still
creates the three legacy scale parameters; the explicit flag is essential.

\section{Parameters and priors}

\begin{table}[h]
\centering\small
\begin{adjustbox}{max width=\linewidth}
\begin{tabular}{llll}
\toprule
Parameter & Prior & Unit & Meaning \\
\midrule
$M_\star$ & log-uniform $[10^6, 10^7]$ & $\Msun$ & stellar mass \\
$M_{\rm rem}$ & log-uniform $[10^4, 3\times10^6]$ & $\Msun$ & remnant mass \\
$a_{\rm rem}$ & log-uniform $[0.3, 20]$ & pc & remnant scale radius \\
$M_\bullet$ & log-uniform $[10^2, 3\times10^5]$ & $\Msun$ & central point mass \\
$\beta_0$ & uniform $[-1, -1/2]$ & --- & central anisotropy (An \& Evans) \\
$\beta_\infty$ & uniform $[-1, 1]$ & --- & outer anisotropy \\
$r_\beta$ & log-uniform $[0.5, 100]$ & pc & anisotropy scale \\
$M_{\rm DM}(<100\,\mathrm{pc})$ & log-uniform $[10^3, 3\times10^7]$ & $\Msun$ & K2 only \\
$r_s$ & log-uniform $[3, 1000]$ & pc & K2 only \\
$s_{\rm inst}$ & uniform, $\sim[0.7,1.3]$ & --- & optional, one per instrument \\
$D$ & normal $(5.43, 0.05)$ & kpc & free in the adopted runs \\
\bottomrule
\end{tabular}
\end{adjustbox}
\caption{Available parameters. Rung 0 fixes anisotropy and removes instrument scales,
leaving 5/7 parameters for K1/K2, including distance. Constant anisotropy gives 6/8;
the implemented varying form gives 8/10. MUSE's diagnostic scale prior is $[0.85,1.15]$;
the Gaia priors are $[0.7,1.3]$.}
\end{table}

\section{Inference}

Nested sampling with \texttt{ultranest}: 400 live points, $\mathrm{d}\ln Z<0.5$, MLFriends by
default with a slice stepper available for the higher-dimensional variants. Posterior mass profiles and residuals are the main scientific products; evidence
compares the stated families after their adequacy and input comparability are checked.

Each run writes an equally weighted posterior, a summary with evidence and per-dataset
$\chi^2$ at the best sampled point, a subset of radial mass profiles, and \texttt{run.yaml}.
The quoted evidence errors are sampler estimates; they do not include uncertainty from
model inadequacy or missing data covariance.

New runs save the actual likelihood arrays and resolved model before sampling: MGE
coefficients, priors, fixed parameters, distance, halo taper and backend, together with
sampler settings, source hashes, git state and package versions. Reports verify the data
snapshot and reproduce the saved best-sample likelihood before plotting. The driver
refuses an existing run directory, including an interrupted run.

Older runs retain their original products. Reports reconstruct them from recorded options
and current data, with a warning and a likelihood check. New evidence comparisons use
the saved observation arrays; legacy comparisons retain conservative input hashes and
the full mock provenance. See \path{docs/CODE_ANALYSIS_AUDIT.md} for the audit and repair
record. These snapshots preserve resolved fit inputs, not a complete software installation
or the raw-catalogue reduction history.

\section{Cross-checks}
\label{sec:crosschecks}

The inference is not allowed to rest on one implementation.

\paragraph{A second Jeans engine.} JamPy's spherical solver runs behind the same likelihood.
The mass components are expanded as projected MGEs so it can consume them, and the anisotropy
family matches JamPy's logistic parametrisation. The non-Gaussian mass-to-MGE conversion
and independent quadrature make this a numerical cross-check, not an exact identity.
At the K1 and cored best samples, total log likelihoods differ by about 0.7, with larger
opposing HST contributions. The cause of that difference has not been isolated.

\paragraph{A distribution-function model.} An AGAMA Cuddeford--Osipkov--Merritt DF, again
behind the same likelihood. It is a \emph{limited} cross-check and is described as such: a
successful construction does not certify a non-negative distribution function, because AGAMA
clips negative values to zero, which can change the realised tracer density (a ratio of
1.045 at 0.11 pc in one trial). Two further limits: its Cuddeford family cannot represent
the anisotropy profile we fit, and the backend softens the point mass with a $10^{-4}$\,pc
Plummer radius, so its central limit is not Keplerian. It is available as a separate diagnostic; there is no automatic DF-positivity veto on
posterior draws. Its CLI family has different anisotropy parameters and cannot replay an
arbitrary Jeans posterior vector.

\paragraph{Published analyses.} Presets reproduce the assumptions of individual papers and
compare conditional results. They are approximations to published assumptions, not
independent reproductions of each paper's full data selection and dynamical method.

\paragraph{Injection.} Earlier K1 mock fits are exploratory checks, not validation of the revised 89-point
ladder. The generator draws symmetric Gaussian errors using the mean of each asymmetric
error pair and clips low values. It is not an exact draw from the fitted split-normal
likelihood. New mock runs record that noise prescription and restore the source run's
model, including its ladder switches. The scientific design of a rung-specific recovery
programme can therefore be chosen explicitly.

\section{What the model does not include}

Stated plainly, because each is a route by which a spurious halo could appear.

\begin{enumerate}
\item \textbf{Sphericity.} \ocen\ is flattened. A spherical model absorbs flattening into
      mass and anisotropy.
\item \textbf{Rotation is removed, not modelled.} The streaming term is subtracted from the
      model's second moment. A rotating model would be preferable.
\item \textbf{Equilibrium.} The outer bins are near plausible pericentric tidal scales.
      Their ratio to $r_J$ depends on the orbit and tidal-radius convention; the data are
      not demonstrably all inside an equilibrium region.
\item \textbf{Diagonal likelihood.} The error-model choice, the rotation-curve uncertainty and
      the shared field template can act coherently across bins. Rotation uncertainty enters
      as per-bin errors, while the raw/eta choice is tested through separate fits. A shared nuisance parameter is the correct treatment; comparing separate
      fits is the chosen alternative.
\item \textbf{A single tracer population.} Mass segregation makes the dispersion depend on
      stellar mass, and the selections differ in magnitude between instruments.
\end{enumerate}

\section{The runs}

The default data contain 89 points. Rung 0 fixes $\beta=0$ and removes scales.
\begin{center}
\begin{adjustbox}{max width=\linewidth}
\begin{tabular}{lrr}
\toprule
Family & $\ln Z$ & $\chi^2$ / number of points\\
\midrule
K1 & $-190.02\pm0.48$ & 772 / 89\\
K2 cored & $-188.94\pm0.36$ & 761 / 89\\
K2 NFW & $-190.09\pm0.35$ & 763 / 89\\
\bottomrule
\end{tabular}
\end{adjustbox}
\end{center}
Rung 0 completes the intended baseline stage. Its weak cored preference and structured
outer residuals provide a reference for the next controlled comparisons. The central fit
is acceptable under the adopted errors. The component ratios motivate varying anisotropy;
they do not directly measure intrinsic $\beta(r)$.

The next stage is constant $\beta$, followed by a specified turnover extension if justified.
The existing monotonic form is available as an intermediate control. Instrument scales are
outside the compared sequence. No rung-1 or turnover posterior on these data is complete.
See \texttt{docs/MODELLING\_PLAN.md} for the authoritative recipes and
\texttt{docs/RUNG0\_ANALYSIS.md} for the conditional numerical results.

From the repository root with \texttt{PYTHONPATH=src}, use
\texttt{python -m ocen\_dm.cli}, or the installed \texttt{ocen} command. The baseline
requires \texttt{fit --family K1 --isotropic --no-scales}; choose a fresh
\texttt{--label} to preserve old results. A bare \texttt{fit} selects the legacy varying
model with scales, not the baseline.

\texttt{--gaia-errors eta} and \texttt{--gaia-rotation ours} are sensitivity variants for
Gaia only. HST always uses the published rotation curve in the current loader. There is
no HST rotation switch; its locally corrected means cannot replace absolute streaming.
Reports preserve these choices and check the saved likelihood. Mock generation from a
source run restores its model independently of the fitted family's stage. Fits are
launched only when explicitly requested.

\begin{thebibliography}{9}
\bibitem[An \& Evans(2006)]{anevans2006} An, J.~H., Evans, N.~W. 2006, ApJ, 642, 752.
  \url{https://ui.adsabs.harvard.edu/abs/2006ApJ...642..752A}
\bibitem[Binney \& Mamon(1982)]{binneymamon1982} Binney, J., Mamon, G.~A. 1982, MNRAS, 200,
  361. \url{https://ui.adsabs.harvard.edu/abs/1982MNRAS.200..361B}
\bibitem[Vasiliev \& Baumgardt(2021)]{vasiliev2021} Vasiliev, E., Baumgardt, H. 2021, MNRAS,
  505, 5978. \url{https://arxiv.org/abs/2102.09568}
\end{thebibliography}

\end{document}
